% =====================================================================
%  5.19 CONCLUSIONES
%  Rubrica: hay que DIFERENCIAR CLARAMENTE las conclusiones tecnicas
%  de las personales. Mezclarlas baja la nota de forma explicita.
%  Por eso van en dos secciones separadas.
% =====================================================================

\chapter{Conclusiones}
\label{cap:conclusiones}

\section{Conclusiones técnicas}

Tres decisiones de diseño facilitaron el desarrollo documentado. La primera es
mantener las reglas de decisión fuera del código: el motor aplica
\texttt{criteria.json} en lugar de codificar STOPP/START v3 como
ramas de programa. Esa separación, descrita en
el~\cref{cap:arquitectura}, hizo posible auditar el cardiovascular
contra la guía~\cite{omahony2023stopp,sacyl2023stopp} y eliminar
criterios inventados o redundantes sin reescribir el evaluador.

La segunda es usar los tests como especificación ejecutable de la
interpretación realizada, no solo como red de seguridad del
software~\cite{beck2002tdd}. Formular un caso y su resultado esperado antes de
implementar la regla convirtió esa expectativa en un contrato comprobable. El
capítulo de pruebas separa el inventario estático de 985 cláusulas de su
resultado: la suite actual no se ha podido ejecutar en el entorno de cierre y,
por tanto, no se presenta como superada. También recoge las amenazas derivadas
de compartir interpretación entre regla y prueba y de no disponer de revisión
independiente completa.

La tercera es la arquitectura de página única sin servidor. En el
contexto de un TFG que maneja datos de salud, evaluar y persistir el
caso solo en el navegador evita un backend de pacientes y reduce la
exposición asociada a ese componente, aunque no elimina los riesgos del
almacenamiento local. El coste, aceptado en
el~\cref{cap:objetivos}, es no cubrir escenarios multiusuario ni
despliegue institucional.

Los límites técnicos son igual de concretos. El catálogo de fármacos y
diagnósticos está acotado a lo modelado: no cubre el arsenal terapéutico
completo. El modelo de datos admite edad y sexo, pero la interfaz no
ofrece un paso de ficha demográfica: esos campos solo llegan por
importación JSON. La analítica sí se captura en la pestaña «Otros» del paso de
diagnósticos. Así, la evaluación automática desde la UI se limita a
medicación, diagnósticos y esos valores; la importación JSON añade los datos
demográficos, aunque el catálogo actual solo consulta la edad; las condiciones
de dosis y duración exigen confirmación manual; y los criterios dependientes
de otros datos quedan sin evaluar. El recuento del catálogo (218 subcriterios
en \texttt{criteria.json}) no coincide 1:1 con los
190 enunciados de la guía: algunos se desglosan en varias reglas y no existe
una matriz que demuestre la correspondencia completa. Por último, las
dependencias diagnósticas y las familias de variantes se modelaron con mayor
detalle en cardiovascular e hipertensión; el resto de sistemas no tiene el
mismo grado de especialización.

\section{Conclusiones personales}

El trabajo se desarrolló con dos perfiles de tutoría distintos: el
tutor del área informática y la cotutora clínica. Esa dualidad obligó a
traducir de forma continua entre restricciones de software (alcance,
tests, interfaz) y matices de la guía. Cuando una duda clínica no
bloqueaba el desarrollo, se documentó y se siguió con la parte técnica;
cuando una petición de interfaz amenazaba el calendario, se dejó en
cuarentena en lugar de fusionarla. Ejemplos de ese filtro son el
historial de casos ---implementado y después eliminado--- y el selector
jerárquico de variantes, del que solo se incorporó la familia de
hipertensión.

Haría dos cosas de otro modo. La primera: registrar las horas desde el
primer día. La dedicación se ha reconstruido a posteriori a partir del
histórico de Git y de las rondas de trabajo, y esa reconstrucción es
menos precisa que un registro contemporáneo. La segunda: anteponer el
paso de datos del paciente (edad, sexo y constantes) a refinamientos de
la taxonomía. Varios criterios del motor ya dependían de esos campos
mientras la UI aún no permitía introducirlos; corregir ese desajuste
antes habría alineado el modelo, las pruebas y el flujo de entrada.

En conjunto, el proyecto muestra la importancia de acotar el alcance, hacer
comprobables las reglas y declarar de forma explícita lo que la evidencia
disponible todavía no permite sostener. Sin pacientes, revisión clínica
independiente completa ni evaluación formal de usabilidad, no se concluye
utilidad clínica ni mejora del trabajo en consulta.
